\documentclass[10pt]{article}
\usepackage[a4paper,margin=2cm]{geometry}
\usepackage{amsmath}
\usepackage{booktabs}
\usepackage{array}
\usepackage{xcolor}
\usepackage{colortbl}
\usepackage{titlesec}
\usepackage{enumitem}
\usepackage{hyperref}
\usepackage{graphicx}
\hypersetup{colorlinks=true,linkcolor=blue!50!black,urlcolor=blue!50!black}
\titlespacing*{\section}{0pt}{1.2ex}{0.6ex}
\titlespacing*{\subsection}{0pt}{0.9ex}{0.4ex}
\setlist{nosep,leftmargin=1.4em}
\definecolor{hdr}{RGB}{235,235,240}
\definecolor{warn}{RGB}{255,244,222}
\newcommand{\cmd}[1]{\texttt{\small #1}}

\title{\vspace{-1.5cm}\textbf{ASearcher: dataset inspection report}\\[2pt]
\large Wave-0 M7 (rung 3, long-horizon web search), CA-Benchmark project}
\author{}
\date{2026-07-25}

\begin{document}
\maketitle
\vspace{-0.6cm}

\noindent\small\textit{Scope.} Direct inspection of the ASearcher datasets (inclusionAI) for rung~3 of
the credit-assignment benchmark, answering the three questions posed: \textbf{how many turns},
\textbf{whether expert trajectories exist}, and \textbf{whether there is a train/eval split}. All counts
come from downloading and parsing the raw files (the HF dataset viewer is broken for these repos);
reproduce with \cmd{scripts/asearcher/inspect\_dataset.py}, assets in \cmd{assets/}. Paper facts are from
arXiv 2508.07976 (verified).\normalsize

\section{Headline answers}
\begin{enumerate}
\item \textbf{Turns: none in the data; horizon is a training-time setting.} ASearcher's data is
(question, answer) pairs with \emph{no} interaction traces, so it carries no turn counts. The paper sets
the RL \emph{turn budget} to \textbf{32 turns} for the 7B/14B agents and \textbf{128 turns} for the
QwQ-32B agent, with emergent behaviour exceeding \textbf{40 tool calls / 150k tokens} per trajectory. Our
benchmark caps rung~3 at \textbf{30 turns} by design, so we sit at the 7B/14B operating point, not the
128-turn regime the ``Beyond Ten Turns'' title advertises.
\item \textbf{Expert trajectories: NO.} Neither training split contains any thought / search-query /
retrieved-doc / turn field --- verified by scanning every row's schema. There is no published SFT /
cold-start / demonstration dataset either. ASearcher trains RL (GRPO) \emph{directly} from base models
with no behaviour-cloning warm start. The closest thing to ``expert'' guidance anywhere is a human
\emph{reference-solution outline} attached to two \emph{eval} sets (GAIA \cmd{Annotator\_Metadata.Steps},
xbench \cmd{reference\_steps}) --- prose plans, not agent rollouts, and eval-only.
\item \textbf{Train/eval split: yes, as two separate repositories.} \cmd{ASearcher-train-data} holds two
\emph{training} splits ($\sim$70.6k QA pairs total, no held-out test inside it); evaluation is a distinct
repo \cmd{ASearcher-test-data} with 10 benchmark subsets. Train and eval never live in one file.
\end{enumerate}

\noindent\textbf{Consequence for our plan.} The proposal already anticipated this: rung~3's SFT
$\pi_\text{base}$ must be \emph{self-generated} (strong-model rollouts in our capped 30-turn setup, or
DR-Tulu remap) because ASearcher ships none. And the solvability diagnostic (M7 step~4) is unavoidable ---
the data gives us questions, but whether \emph{our} 4B agent under \emph{our} wiki-18 corpus at 30 turns
can answer them is exactly what the go/no-go gate must measure.

\section{The training data: \cmd{inclusionAI/ASearcher-train-data}}
Two JSONL files, both \textbf{QA pairs only}, Apache-2.0. Data is \emph{synthesized}: a prompt-based LLM
agent starts from 14,107 seed questions and applies ``Injection'' (add related facts) and ``Fuzz''
(obscure identifying details) operations, then a 3-step quality/difficulty/uniqueness filter. It is
QA-pair synthesis, \emph{not} trajectory synthesis.

\begin{center}\small
\begin{tabular}{lrl}
\toprule
Split & Rows & Fields \\
\midrule
\cmd{ASearcher-Base-35k} & 35{,}583 & \cmd{question}, \cmd{answer} (list), \cmd{source}, \cmd{aug\_answer} (aliases), \cmd{qid} \\
\cmd{ASearcher-LRM-35k} & 35{,}054 & \cmd{question}, \cmd{answer} (string), \cmd{source}, \cmd{id}, \cmd{idx} \\
\bottomrule
\end{tabular}
\end{center}

\noindent The two splits are genuinely different products, not a random partition:
\begin{itemize}
\item \textbf{Base-35k} --- shorter multi-hop questions (median \textbf{23} words, p95 60), a list-valued
answer, and \emph{heavy alias augmentation} (\cmd{aug\_answer}: median 4 paraphrases, max 1454). Sources:
\cmd{compose\_chain\_qa} 8{,}537, \cmd{opensource-fm8k} 7{,}991, \cmd{opensource} 7{,}927,
\cmd{compose\_group\_qa} 6{,}375, and \cmd{compose\_chain\_qa\_invalid} \textbf{4{,}753}. Note that last
tag: $\sim$13\% of Base is flagged \emph{invalid} yet shipped --- we must decide to drop or keep it (drop,
unless the diagnostic says otherwise).
\item \textbf{LRM-35k} --- longer, more convoluted ``chain'' questions (median \textbf{43} words, p95 88,
max 639) aimed at large-reasoning-model agents; string answer, no aliases. Sources: 24{,}026 rows carry
\emph{no source tag} (bare numeric ids), \cmd{chain\_qa\_fuzzy} 3{,}527, \cmd{math} 2{,}930,
\cmd{train\_v1:webwalker\_hard} \textbf{2{,}572}, \cmd{train\_v1:compose\_group\_qa} 1{,}999.
\end{itemize}

\noindent The question-length gap (23 vs 43 median words) is itself a within-rung difficulty axis. The
\cmd{math} slice in LRM and the missing aliases mean LRM is the harder, noisier half --- relevant to which
split we draw the frozen training set from.

\section{The eval data: \cmd{inclusionAI/ASearcher-test-data}}
A separate repo, one \cmd{test} split per benchmark subdirectory, all QA pairs (+ metadata), no
trajectories. Row counts verified by download:

\begin{center}\small
\begin{tabular}{lrl}
\toprule
Subset & Rows & Notes \\
\midrule
GAIA & 103 & has human \cmd{Annotator\_Metadata.Steps} (reference plan); Level 1--3 \\
frames (FRAMES) & 824 & multi-constraint; gold Wikipedia link chains + \cmd{reasoning\_types} \\
xbench-deepsearch & 100 & Chinese deep-search; has \cmd{reference\_steps}; canary-guarded \\
Bamboogle & 125 & short multi-hop QA \\
2WikiMultihopQA\_rand1000 & 1{,}000 & random-1000 subsample \\
HotpotQA\_rand1000 & 1{,}000 & random-1000 subsample \\
Musique\_rand1000 & 1{,}000 & random-1000 subsample \\
NQ\_rand1000 & 1{,}000 & random-1000 subsample \\
PopQA\_rand1000 & 1{,}000 & random-1000 subsample \\
TriviaQA\_rand1000 & 1{,}000 & random-1000 subsample \\
\bottomrule
\end{tabular}
\end{center}
\noindent Total 7{,}277 eval items across 10 benchmarks. The paper's headline (web setting, ASearcher-Web-QwQ)
reports GAIA $\sim$53--59, xbench $\sim$42--51, FRAMES $\sim$71--75 (exact numbers are checkpoint-dependent
between the arXiv version and the repo leaderboard --- pin the checkpoint before citing).

\section{Evaluation metric (and the training reward --- they differ)}
ASearcher uses \textbf{two different scoring functions} depending on whether it is training or evaluating;
this matters because our proposal's ``EM, verifiable'' for rung~3 matches only the \emph{training} one.

\smallskip
\noindent\textbf{Eval metric = LLM-as-Judge (``MBE'', the headline reported numbers).} From
\cmd{evaluation/evaluate.py}: the predicted answer and the gold answer are handed to a judge LLM with a
binary-equivalence prompt --- \emph{``determine if the predicted answer is equivalent to the labeled
answer $\ldots$ respond Correct / Incorrect''} --- returning a rationale plus a \cmd{judgement}; MBE $=1$
iff \cmd{judgement=correct}, averaged over the set.
\begin{itemize}
\item \textbf{Judge model:} \cmd{Qwen2.5-72B-Instruct} (self-hosted default), or \cmd{gpt-4o-mini} via
the OpenAI path.
\item \textbf{Aggregation:} reported as \textbf{Avg@k} (mean correctness over $k$ sampled rollouts) and
\textbf{Pass@k} (correct if any of $k$), eval temperature 0.6, $k=$ \cmd{--pass-at-k} (docs advise raising
$k$ for stability). This judge metric is what they use on the free-form sets (GAIA, xbench, FRAMES) where
exact match under-counts.
\end{itemize}

\noindent\textbf{Training reward = rule-based, verifiable, no judge in the loop} (\cmd{ASearcher/utils/rewards.py}):
extract the last \cmd{<answer>$\ldots$</answer>}, then score by \textbf{EM} with SQuAD-style normalization
(lowercase, strip articles, drop punctuation, whitespace-fix, plus \cmd{True}$\to$\cmd{yes} mapping and
Chinese-token handling); \textbf{sub-EM} (substring) and token-\textbf{F1} variants are also implemented.
No extractable answer $\Rightarrow$ reward 0. The paper stresses there is deliberately no external LLM in
the agent/training loop.

\smallskip
\noindent\textbf{For us:} rung-3 RL reward = their training EM/sub-EM/F1 (we already have this from
rung~1's Search-R1 harness). For the transfer eval on GAIA/xbench/FRAMES we adopt the judge track ---
with our pinned \cmd{gpt-oss-120b} substituted for Qwen2.5-72B --- which is exactly the uniform-judge track
in the harness (§6). \textbf{Caveat:} any ASearcher SOTA number we cite must record \emph{which judge} and
\emph{which $k$}, or it is not comparable (different judge model $\to$ different MBE).

\section{Tool-calling interface (and what changes offline)}
\textbf{The action grammar is identical online and offline} --- only the backend behind each tool changes.
The full ASearcher agent (\cmd{agent/asearcher.py}, \cmd{asearcher\_reasoning.py}) emits three tags,
generation stopping at each closing tag (\cmd{</search>}, \cmd{</access>}, \cmd{</answer>}):
\begin{center}\small
\begin{tabular}{lll}
\toprule
Tool tag & meaning & backend: web setting $\to$ local (offline) setting \\
\midrule
\cmd{<search>query</search>} & retrieval & Serper API $\to$ local dense-retriever \cmd{/retrieve} (e.g.\ e5 over wiki-18) \\
\cmd{<access>url</access>} & browse a page & Jina live extraction $\to$ local \cmd{/access} over a \emph{prebuilt} url$\to$page store \\
\cmd{<answer>$\ldots$</answer>} & terminate & (returns final answer for scoring) \\
\bottomrule
\end{tabular}
\end{center}

\noindent\textbf{Does offline call the same tools? Yes --- both \cmd{<search>} and \cmd{<access>} are
supported offline.} The bundled \cmd{tools/local\_retrieval\_server.py} exposes \emph{both} a
\cmd{POST /retrieve} (dense search over the corpus) and a \cmd{POST /access} endpoint. The crucial
difference: \cmd{/access} is backed by a \cmd{PageAccess} class that serves pages from a \emph{pre-built
file keyed by URL} (Wikipedia pages, with \cmd{index.php/}$\to$\cmd{index.php?title=} normalization) --- so
offline \cmd{<access>} can only open URLs that exist in that store, whereas online Jina fetches
\emph{any} URL live. \cmd{<search>} works over any corpus with a dense index (our existing wiki-18/e5).

\smallskip
\noindent\textbf{Tool-calling options (three agent variants ship in the repo):}
\begin{itemize}
\item \cmd{asearcher} --- 2 tools: \cmd{<search>} + \cmd{<access>} + \cmd{<answer>}. The design that
produces the long-horizon behaviour (search then browse specific pages).
\item \cmd{asearcher-reasoning} --- same 2 tools wrapped with an explicit \cmd{<thought>} block (for
reasoning models).
\item \cmd{search\_r1} (\cmd{agent/search\_r1.py}) --- \textbf{search-only}: \cmd{<search>} + \cmd{<answer>}
(Search-R1 grammar), \emph{no} \cmd{<access>}. This is the variant that matches our existing rung-1 infra.
\end{itemize}

\noindent\textbf{Decision this forces for rung~3 (flag for the manifest).} Our current retrieval stack is
\emph{search-only} (rung~1 serves \cmd{/retrieve} over wiki-18, no page store). Two options:
\textbf{(a) search-only rung~3} using the \cmd{search\_r1} agent --- zero new infra, reuses rung~1 exactly,
but drops the browse tool that is central to ASearcher's ``deep search'' and shortens realized horizons;
\textbf{(b) full 2-tool} --- additionally build the wiki-18 url$\to$page store and stand up \cmd{/access},
reproducing ASearcher-local faithfully at real extra W0 cost. This choice interacts with the horizon claim
(RQ2): without \cmd{<access>}, rung~3's practical turn count and its separation from rung~1 both shrink, so
\textbf{(b) is the more faithful long-horizon rung} if the budget allows. Recommend deciding at the M7 gate
alongside the solvability diagnostic, since answerability under search-only vs search+browse can differ
sharply.

\section{Two leakage vectors we must handle (important)}
\begin{enumerate}
\item \textbf{ASearcher train $\cap$ WebWalkerQA.} The LRM split contains \textbf{2{,}572
\cmd{webwalker\_hard} items}. WebWalkerQA is one of \emph{our} rung-3 transfer-eval benchmarks (proposal
§4). Training on webwalker-derived questions and then reporting WebWalkerQA transfer would be leakage ---
the proposal's ``fuzzy dedup vs WebWalkerQA'' step is therefore load-bearing, and may require dropping the
whole \cmd{webwalker\_hard} slice, not just exact-match hits.
\item \textbf{ASearcher eval $\cap$ our rung-1 SearchQA.} \cmd{ASearcher-test-data} \emph{is}
HotpotQA / 2Wiki / Musique / NQ / TriviaQA / PopQA / Bamboogle --- the same families we \emph{train on and
eval} at rung~1. We must not let rung-3 training questions (composed from these) overlap rung-1 eval, and
we should use \emph{our own} rung-1 eval splits, not ASearcher's rand1000 subsamples, to keep cross-rung
comparisons clean. (Spot check: 0/125 exact-normalized overlap between ASearcher-train and the Bamboogle
test file --- reassuring for exact matches, but the composed questions demand fuzzy dedup, not exact.)
\end{enumerate}

\section{Environment / answerability note}
ASearcher ships \emph{two} settings: a \textbf{local-RAG} setting (Wikipedia-2018 corpus --- the same
wiki-18 dump we already host for rung~1) for its Local models, and a \textbf{live-web} setting (Serper +
Jina) for its Web models. Its strongest published numbers are the \emph{live-web} QwQ-32B agent. We run the
\textbf{local-RAG, 30-turn, Qwen3-4B} configuration --- strictly harder on answerability than live web and
much smaller than QwQ-32B. \textbf{This is precisely why the M7 solvability diagnostic gates the rung:} the
questions were synthesized to be hard for a 7B--32B agent with strong retrieval; our answerable-fraction
check on 200--300 items under wiki-18 decides regime (i)/(ii)/(iii) before we spend training budget.

\section{Challenges (rung-3 specific)}
\begin{itemize}
\item \textbf{No warm start.} No expert trajectories $\Rightarrow$ rung-3 SFT $\pi_\text{base}$ must be
generated by us (strong-model rollouts in-setup, EM-verified), which is extra W0 inference cost and its own
quality-control pass.
\item \textbf{Answerability is unknown and setting-dependent.} Data is tuned for live-web / large models;
under our capped local-RAG 4B agent the learnable-band fraction is the open question --- the go/no-go gate.
\item \textbf{Leakage on both sides} (webwalker into training, rand1000 benchmarks overlapping rung~1) ---
fuzzy dedup required, not exact.
\item \textbf{Long, variable horizon} (paper 32--128 turn budgets; we cap 30) with retrieved-doc context
that will not fit full-history --- the rung-3 truncation policy (per-doc caps + drop-oldest
\cmd{<information>}) is binding here, unlike the QA rung.
\item \textbf{Noisy/heterogeneous training pool} --- a 13\% ``invalid'' slice in Base, a 40\% source-less
slice and a \cmd{math} slice in LRM; the frozen training draw needs an explicit inclusion policy.
\end{itemize}

\section*{Reproduce}
\cmd{python scripts/asearcher/inspect\_dataset.py --out <dir>} downloads both repos (public, no token) and
regenerates \cmd{train\_stats.json}, \cmd{test\_stats.json}, \cmd{sample\_rows.json}, \cmd{summary.md}.

\end{document}
